% Generated draft: requires mathematical and source review.
\documentclass[11pt]{article}
\usepackage[margin=25mm]{geometry}
\usepackage{fontspec}
\setmainfont{Latin Modern Roman}
\usepackage{amsmath,amssymb,mathtools,unicode-math}
\setmathfont{Latin Modern Math}
\usepackage{enumitem,needspace,xurl,hyperref}
\setlength{\parindent}{0pt}
\setlength{\parskip}{5pt}
\setlength{\emergencystretch}{3em}
\newcommand{\N}{\mathbb{N}}
\newcommand{\Z}{\mathbb{Z}}
\newcommand{\Q}{\mathbb{Q}}
\newcommand{\R}{\mathbb{R}}
\newcommand{\C}{\mathbb{C}}
\newcommand{\F}{\mathbb{F}}
\newcommand{\A}{\mathbb{A}}
\newcommand{\PP}{\mathbb{P}}
\newcommand{\E}{\mathbb{E}}
\newcommand{\eps}{\varepsilon}
\newcommand{\dd}{\,\mathrm{d}}
\begin{document}
\section*{UnsolvedMath ID: 701519}
\textbf{Source problem code:} catalogue entry
\subsection{Definitions and mathematical statement}
The supplied record does not explicitly specify the allowed values of the parameters \(m,n\). The formulation below uses positive integers; this intended parameter range requires confirmation.

For \(m,n\in\Z_{\geq 1}\), let
\[
Y(m,n)=\left\langle a_1,\ldots,a_m\ \middle|\ a_k^n=1\ (1\leq k\leq m),\quad (a_k^i a_l^i)^2=1\ \left(1\leq k<l\leq m,\ 1\leq i\leq\left\lfloor\frac n2\right\rfloor\right)\right\rangle.
\]
All indices in this presentation are integers, and \(1\) denotes the group identity. An empty range contributes no relations. Subject to confirmation of the parameter range, the requested assertion is
\[
\forall m,n\in\Z_{\geq 1},\qquad |Y(m,n)|<\infty.
\]

The following are contextual claims supplied with the problem, not independently verified assertions of this draft. Write \(C_n\) for the cyclic group of order \(n\), \(A_r\) for the alternating group of degree \(r\), and \(C_2^m\) for the direct product of \(m\) copies of \(C_2\). The source states
\[
Y(m,2)\cong C_2^m,\qquad Y(m,3)\cong A_{m+2},
\]
attributes the latter presentation to Carmichael, and states that \(Y(m,4)\) has order \(2^{m(m+3)/2}\), nilpotency class \(3\), and exponent \(4\). The class claim needs a qualification on \(m\): at \(m=1\), the defining presentation gives the cyclic group \(C_4\), which has nilpotency class \(1\).

Let \(\F_2[C_n]\) be the group algebra and let its augmentation ideal be
\[
\omega(C_n)=\ker\!\left(\F_2[C_n]\longrightarrow\F_2,\quad \sum_{g\in C_n}c_g g\longmapsto\sum_{g\in C_n}c_g\right).
\]
The source describes \(Y(2,n)\) as the natural extension of \(\omega(C_n)\) by \(C_n\), attributing its presentation to Coxeter, and states \(Y(3,n)\cong SL(2,\omega(C_n))\). The extension data and the meaning of \(SL(2,\omega(C_n))\), whose ring argument is generally nonunital, are not specified.

For odd \(n\), the source additionally says that \(Y(m,n)\) has the following presentation, but labels the presented group with a lowercase \(y\):
\[
\begin{aligned}
y(m,n)=\big\langle a,\mathbb S_m\ \big|\;&a^n=1,\\
&[\tau_{12}^{a^i},\tau_{12}]=1\quad\left(1\leq i\leq\left\lfloor\frac n2\right\rfloor\right),\\
&\tau_{12}^{1+a+\cdots+a^{n-1}}=1,\\
&\tau_{i,i+1}a\tau_{i,i+1}=a^{-1}\quad(2\leq i\leq m-1)\big\rangle.
\end{aligned}
\]
Here \(\mathbb S_m\) is the symmetric group on \(m\) letters, with its own group relations understood, and \(\tau_{ij}\) denotes the transposition interchanging \(i\) and \(j\). This display requires \(m\geq2\). The source does not explain the relationship between its uppercase and lowercase group names, its conjugation convention, or the word represented by the formal-sum exponent \(1+a+\cdots+a^{n-1}\). Consequently this auxiliary presentation is retained as supplied context and is not used as a verified equivalent definition of \(Y(m,n)\).

The source reports the following examples, with its classical-group notation retained:
\[
\begin{aligned}
y(3,5)&\cong SL(2,16)\cong\Omega^{-}(4,4),\\
y(4,5)&\cong Sp(2,16)\cong\Omega(5,4),\\
y(5,5)&\cong SU(4,16)\cong\Omega^{-}(6,4),\\
y(6,5)&\cong 4^6\Omega^{-}(6,4),\\
y(3,7)&\cong SL(2,8)^2\cong\Omega^{+}(4,8),\\
y(4,7)&\cong Sp(4,8)\cong\Omega(5,8).
\end{aligned}
\]
The classical-group parameter conventions and the extension notation \(4^6\Omega^{-}(6,4)\) require clarification before these examples can be used precisely. The source describes connections with classical groups in characteristic \(2\), attributes extensive computations confirming this trend to Felsch, Neub\"user and O'Brien, and mentions Bott periodicity and Clifford algebras without giving a precise additional assertion.
\subsection{Short English statement}
S. Sidki's problem asks whether every group \(Y(m,n)\) defined by the displayed generators and relations is finite. The intended range of \(m,n\) is unstated; positive integers are used provisionally. The supplied registry labels the problem open. Its auxiliary presentation and contextual examples contain unresolved notation and qualifications and require human mathematical review.
\subsection{Sources}
\begin{itemize}
\item Supplied local UnsolvedMath registry record: permanent ID \(701519\), exact source problem code \texttt{catalogue entry}, title \emph{Kourovka Notebook Problem 18.90}. The record labels the problem \emph{open}; this is supplied registry status, not an independently verified literature assessment. Its recorded update date is \texttt{2026-10-16T14:00:00Z}, and its upstream revision is \texttt{d371ee917ac9647fbb06cd6fc529deec6f3dda82}.
\item Supplied primary-source attribution: \emph{The Kourovka Notebook (Unsolved Problems in Group Theory)}, Problem 18.90, proposed by S. Sidki in the 18th issue (2014). The background also identifies the 21st issue edition (2026), \texttt{arXiv:1401.0300v48}. The supplied collection description names E. I. Khukhro and V. D. Mazurov as editors. These bibliographic details are reproduced from the record; the cited edition has not been consulted for this draft.
\item Contextual attributions supplied in the statement: Carmichael for the presentation associated with \(Y(m,3)\); Coxeter for the presentation associated with \(Y(2,n)\); and Felsch, Neub\"user and O'Brien for extensive computations. No separate bibliographic citations or links for these contributions are supplied.
\end{itemize}
\end{document}
